\documentclass[a4paper,11pt]{article}
\usepackage{exam-style-842-unit}
\usepackage{needspace}
\examinehead{842 信息化概论 · 第06单元练习册}{练习一\quad 物联网概念与体系架构}
\begin{document}

\unitcover{842 信息化概论 · 单元练习册}{第06单元\quad 物联网与传感技术}{练习一\quad 物联网概念与体系架构（共 100 分）}{约 70 分钟}{物联网的定义与两层含义；基本特征；教材四层体系架构（感知层·网络层·支撑层·应用层）及各层功能与设备；物联网关键技术（感知·传输·支撑）；与互联网、传感网、泛在网的辨析；发展历程、标准化与前景。}

\begin{knowmap}
\textbf{定义：}物联网（Internet of Things，IoT）是通过射频识别（RFID）、红外感应器、全球定位系统、激光扫描器等信息传感设备，按约定的协议把任何物品与互联网相连接，进行信息交换和通信，以实现智能化识别、定位、跟踪、监控和管理的一种网络\\[2pt]
\textbf{两层含义：}核心和基础仍然是互联网 $\rightarrow$ 用户端延伸和扩展到物品与物品之间，可进行物与物、人与物、人与人的信息交换\\[2pt]
\textbf{三大特征：}互联网特征（物品互联互通）· 识别与通信特征（自动识别、物物通信、全面感知）· 智能处理（云计算、模糊识别等智能计算）\\[2pt]
\textbf{四层架构：}感知层（采集物理事件和数据）$\rightarrow$ 网络层（接入与传输，主要设备是网关）$\rightarrow$ 支撑层（云计算与海量存储的超级平台）$\rightarrow$ 应用层（行业应用平台与管理界面）\\[2pt]
\textbf{关键技术：}感知技术（RFID·传感器·GPS·多媒体采集处理·二维码）· 传输技术（移动通信网·互联网·无线网络·近程通信）· 支撑技术（嵌入式计算·云计算·分布式并行计算·人工智能·数据库与数据挖掘·多媒体与虚拟现实）\\[2pt]
\textbf{发展脉络：}1995 年《未来之路》提及 $\rightarrow$ 1998 年 Auto-ID 中心应用于物品 $\rightarrow$ 1999 年 EPC global 真正提出 $\rightarrow$ 2005 年 ITU 报告拓展到传感网络 $\rightarrow$ 2009 年“智慧地球”“感知中国” $\rightarrow$ 2010 年发展元年 $\rightarrow$ 2012 年八领域国家物联网应用示范工程
\end{knowmap}

\begin{keyterms}
\textbf{物联网}：通过 RFID、红外感应器、GPS 等信息传感设备把物品与互联网相连，实现智能化识别、定位、跟踪、监控和管理的网络。\\
\textbf{感知层}：采集物理世界中发生的物理事件和数据，利用 RFID、传感器、GPS、GIS 等实现物理世界的全面感知。\\
\textbf{网络层}：通过互联网、移动通信网、卫星通信网等对感知层信息进行接入和传输，主要设备是能接入各种异构网络的网关。\\
\textbf{支撑层}：在高性能网络计算环境下把海量信息资源整合成大型智能网络，为上层的服务管理和行业应用提供网络计算超级平台。\\
\textbf{应用层}：物联网系统结构的最高层，面向各类行业建立管理平台和运行平台，提供绿色农业、智能交通等服务。\\
\textbf{传感网}：由大量智能传感器组成的信息感知网络，是物联网的前身，现在是物联网的组成部分。\\
\textbf{泛在网}：无所不在的网络，是物联网的发展远景，物联网是泛在网的前期。\\
\textbf{感知技术}：用于物联网底层感知信息的技术，包括 RFID、传感器、GPS、多媒体采集处理、二维码等。\\
\textbf{传输技术}：汇聚感知数据并实现物联网数据传输的技术，包括移动通信网、互联网、无线网络、卫星通信、短距离无线通信等。\\
\textbf{支撑技术}：用于物联网数据处理和利用的技术，包括嵌入式计算、云计算、分布式并行计算、人工智能、数据库与数据挖掘、多媒体与虚拟现实等。\\
\textbf{二维码技术}：用黑白相间的几何图形按规律记录数据符号信息，能在横向和纵向同时表达信息，在很小面积内表达大量信息。\\
\textbf{RFID}：射频识别技术，物联网中让物品“开口说话”的关键技术，标签上存储规范而具有互用性的信息，可自动采集到中央信息系统。\\
\textbf{智慧地球}：2009 年 IBM 提出的构想，互联网加物联网构成智慧地球，物联网是其不可或缺的重要部分。\\
\textbf{感知中国}：2009 年我国提出的物联网发展战略计划，无锡物联网产业园区是其首个前沿阵地。
\end{keyterms}

\needspace{6cm}
\begin{contrast}{物联网四层架构（每行依次为：主要功能·主要设备或技术·典型举例）}
\citem{感知层}{采集物理世界中发生的物理事件和数据}{RFID、传感器、GPS、GIS、二维码}{电子耳标、土壤温湿度传感器、视频监控}
\citem{网络层}{对来自感知层的信息进行接入和传输}{互联网、移动通信网、卫星通信网、异构网络网关}{4G/5G、光纤宽带、无线局域网}
\citem{支撑层}{把海量信息资源整合成可互联互通的大型智能网络，提供网络计算超级平台}{大型计算机群、海量网络存储设备、云计算设备}{农业云平台、海量数据存储与计算}
\citem{应用层}{面向行业建立实际应用的管理平台和运行平台}{各类用户界面显示设备及其他管理设备}{绿色农业、智能交通、智能家居、远程医疗}
\end{contrast}

\needspace{6cm}
\begin{contrast}{物联网 / 互联网 / 传感网 / 泛在网（每行依次为：定位·相互关系·关键判据）}
\citem{物联网}{物物相连的互联网}{是互联网的延伸和扩展，互联网是其核心和基础}{用户端扩展到物品与物品之间}
\citem{互联网}{全球范围的信息交互网络}{物联网的基础网络，物联网在其基础上延伸}{连接的是计算机等终端}
\citem{传感网}{由大量智能传感器组成的信息感知网络}{物联网的前身，现在是物联网的组成部分}{侧重信息感知与采集}
\citem{泛在网}{无所不在的网络社会}{物联网的发展远景，物联网是泛在网的前期}{强调任何人、任何物、任何时间、任何地点的连接}
\end{contrast}

\needspace{6cm}
\begin{contrast}{业界三层 / 教材四层（每行依次为：层次划分·与另一方案的关系·答题说明）}
\citem{业界三层}{感知层、网络层、应用层}{把计算与存储支撑并入网络层、应用层来理解}{历年真题按三层作答亦应给分}
\citem{教材四层}{感知层、网络层、支撑层、应用层}{在网络层与应用层之间单列支撑层（云计算、海量存储的超级平台）}{本册以教材四层为主，答题时说明二者对应关系}
\end{contrast}

\qtypename{一、单项选择题}{10 题，每题 3 分，共 30 分}
\q 通过射频识别、红外感应器、全球定位系统、激光扫描器等信息传感设备，按约定的协议把任何物品与互联网相连接，以实现智能化识别、定位、跟踪、监控和管理的网络是（\quad）
\choicesline{A. 传感网}{B. 物联网}{C. 泛在网}{D. 互联网}

\q 物联网的核心和基础是（\quad）
\choicesline{A. 互联网}{B. 传感网}{C. 泛在网}{D. 云计算}

\q 按教材第 7 章，物联网体系结构自下而上依次分为（\quad）
\choices{A. 感知层、网络层、应用层}{B. 感知层、网络层、支撑层、应用层}
{C. 网络层、感知层、应用层}{D. 应用层、支撑层、网络层、感知层}

\q 教材物联网四层架构中，为上层服务管理和大规模行业应用提供高效、可靠、可信的网络计算超级平台的是（\quad）
\choicesline{A. 感知层}{B. 网络层}{C. 支撑层}{D. 应用层}

\q 物联网体系结构中，主要设备是能接入各种异构网络的网关的是（\quad）
\choicesline{A. 感知层}{B. 网络层}{C. 支撑层}{D. 应用层}

\q 在物联网中，让物品“开口说话”的关键技术是（\quad）
\choicesline{A. 二维码技术}{B. 射频识别技术}{C. GPS 定位技术}{D. 多媒体采集技术}

\q 利用 RFID、传感器、GPS、GIS 等随时随地获取物体信息、实现物理世界全面感知的是（\quad）
\choicesline{A. 感知层}{B. 网络层}{C. 支撑层}{D. 应用层}

\q 传感网与物联网的关系是（\quad）
\choices{A. 传感网是物联网的前身}{B. 物联网是传感网的前身}
{C. 二者完全等同}{D. 传感网是泛在网的发展远景}

\q 2009 年提出“智慧地球”构想的公司是（\quad）
\choicesline{A. IBM}{B. 英特尔}{C. 微软}{D. 谷歌}

\q 1999 年真正提出物联网概念的是（\quad）
\choices{A. 比尔·盖茨的《未来之路》}{B. MIT Auto-ID 中心的 Kevin Ashton}
{C. EPC global 的 Auto-ID 中心}{D. ITU《ITU 互联网报告 2005：物联网》}

\qtypename{二、多项选择题}{4 题，每题 4 分，共 16 分；每题有两个或两个以上正确项}
\q 物联网的基本特征包括（\quad）
\choicesline{A. 互联网特征}{B. 识别与通信特征}{C. 智能处理}{D. 高能耗传输}

\q 关于物联网与互联网的关系，下列说法正确的有（\quad）
\choices{A. 核心和基础仍是互联网}{B. 是互联网的延伸和扩展}
{C. 用户端扩展到物品之间}{D. 完全独立于互联网}

\q 物联网感知层可以使用的信息获取技术有（\quad）
\choicesline{A. RFID}{B. 传感器}{C. GPS}{D. 异构网络网关}

\q 关于物联网的发展前景，下列说法正确的有（\quad）
\choices{A. EPoSS 预测 2020 年之后实现全智能化}{B. 物物互联业务将达人际通信业务的 30 倍}
{C. 被称为下一个万亿美元级的通信业务}{D. 大规模应用已不存在成本和安全挑战}

\qtypename{三、判断题}{6 题，每题 2 分，共 12 分；正确的画“√”，错误的画“×”并写出改正后的表述}
\q 物联网是物物相连的互联网，其核心和基础仍然是互联网。\quad（\quad）\par\noindent 改正：\fixblank

\q 教材第 7 章把物联网体系结构划分为感知层、网络层、应用层三个层次。\quad（\quad）\par\noindent 改正：\fixblank

\q 支撑层的主要设备包括大型计算机群、海量网络存储设备和云计算设备。\quad（\quad）\par\noindent 改正：\fixblank

\q 网络层主要负责采集物理世界中发生的物理事件和数据。\quad（\quad）\par\noindent 改正：\fixblank

\q 传感网是物联网的前身，现在是物联网的组成部分；泛在网是物联网的发展远景。\quad（\quad）\par\noindent 改正：\fixblank

\q 物联网的用户端只限于人与人之间的信息交换，不涉及物与物之间的通信。\quad（\quad）\par\noindent 改正：\fixblank

\qtypename{四、配对题}{6 对，每对 1 分，共 6 分；把 A 栏与 B 栏对应内容配对}
\q 将下列概念与其对应内容配对，把 B 栏字母填入括号。
\par\vspace{0.2em}
\noindent\begin{tabular}{@{}p{0.44\textwidth}@{\hspace{1.2em}}p{0.44\textwidth}@{}}
（1）感知层 \hfill（\quad） & A. 让物品“开口说话”的关键技术 \\
（2）网络层 \hfill（\quad） & B. 绿色农业、智能交通等行业应用平台 \\
（3）支撑层 \hfill（\quad） & C. 接入各种异构网络的网关 \\
（4）应用层 \hfill（\quad） & D. 采集物理事件和数据 \\
（5）RFID \hfill（\quad） & E. 云计算、海量存储的超级平台 \\
（6）GPS \hfill（\quad） & F. 全球定位系统，采集移动物体信息 \\
\end{tabular}\par\vspace{0.3em}

\newpage
\qtypename{五、简答题}{3 题，每题 8 分，共 24 分}
\q 简述物联网的定义及其两层含义，并说明物联网的基本特征。\anslines{2.4cm}

\q 简述教材中物联网体系结构四个层次的主要功能与主要设备。\anslines{2.4cm}

\q 简述物联网、互联网、传感网、泛在网四者之间的区别与联系。\anslines{2.4cm}

\newpage
\qtypename{六、综合题}{1 题，共 12 分}
\q 某现代农业产业园在稻田部署土壤温湿度传感器、小型气象站和视频监控，为农机安装北斗终端，给稻米包装赋二维码；数据经 4G 和光纤传至农业云平台，平台开发手机应用向农户推送灌溉建议、向消费者提供溯源查询。请回答：（1）按教材第 7 章的四层架构，把材料中的要素分别归入相应层次，并说明支撑层为何不能省略；（2）若按业界通行的三层架构作答，应如何与教材四层对应；（3）说明该园区体现了物联网的哪些基本特征。\anslines{5cm}

\end{document}
